\documentclass[11pt]{article}

% Shared notes settings for Elements of AIML lectures (UPES).
% Keep this file free of \documentclass to allow reuse via \input.

\usepackage[utf8]{inputenc}
\usepackage[T1]{fontenc}
\usepackage[a4paper,margin=1in]{geometry}
\usepackage{hyperref}
\usepackage{xurl}
\usepackage{booktabs}
\usepackage{amsmath}
\usepackage{amssymb}
\usepackage{listings}
\usepackage{xcolor}
\usepackage{enumitem}
\usepackage{parskip}

% Code listings (general-purpose).
\lstset{
  basicstyle=\ttfamily\small,
  keywordstyle=\color{blue},
  commentstyle=\color{gray},
  stringstyle=\color{teal},
  showstringspaces=false,
  columns=fullflexible,
  frame=single,
  framerule=0.2pt,
  breaklines=true
}

\setlist[itemize]{topsep=0.3em,itemsep=0.2em}
\setlist[enumerate]{topsep=0.3em,itemsep=0.2em}

% Simple per-section source line.
\newcommand{\notesources}[1]{\par\noindent{\small\color{gray}\textit{Sources:} \sloppy #1\par}}

% Unit-wise source strings for this subject (used by slides/notes).
% Path is relative to the lecture `latex/` folder (the compilation CWD).
\IfFileExists{../../../_shared/aiml-sources.tex}{%
  % Source strings for Elements of AIML (used by slides + notes).

\newcommand{\AIMLRepoURL}{https://github.com/tali7c/Elements-of-AIML}

\newcommand{\AIMLLabSources}{%
Provost and Fawcett, \textit{Data Science for Business}.;%
McKinney, \textit{Python for Data Analysis}.;%
WEKA Documentation (Waikato Environment for Knowledge Analysis), accessed 2026-02-28.;%
Matplotlib and Seaborn documentation, accessed 2026-02-28.%
}

\newcommand{\AIMLUnitISources}{%
Rich and Knight, \textit{Artificial Intelligence}, McGraw Hill, 2017.;%
Russell and Norvig, \textit{Artificial Intelligence: A Modern Approach} (4e), Ch. 1--2 (Introduction, Intelligent Agents).;%
Poole and Mackworth, \textit{Artificial Intelligence: Foundations of Computational Agents} (2e), selected sections.%
}

\newcommand{\AIMLUnitIISources}{%
Russell and Norvig, \textit{Artificial Intelligence: A Modern Approach} (4e), Ch. 7--9 (Logical Agents, First-Order Logic, Inference).;%
Huth and Ryan, \textit{Logic in Computer Science} (2e), selected sections on propositional and first-order logic.;%
Genesereth and Nilsson, \textit{Logical Foundations of Artificial Intelligence}, selected chapters.%
}

\newcommand{\AIMLUnitIIISources}{%
Mueller and Massaron, \textit{Machine Learning for Dummies}, For Dummies, 2016.;%
Mitchell, \textit{Machine Learning}, selected chapters on learning basics and evaluation.;%
Scikit-learn documentation, accessed 2026-02-28.%
}

\newcommand{\AIMLUnitIVSources}{%
Mueller and Massaron, \textit{Machine Learning for Dummies}, For Dummies, 2016.;%
James, Witten, Hastie, and Tibshirani, \textit{An Introduction to Statistical Learning} (2e), selected chapters.;%
Scikit-learn documentation, accessed 2026-02-28.%
}

\newcommand{\AIMLUnitVSources}{%
NIST, \textit{AI Risk Management Framework (AI RMF 1.0)}, 2023.;%
Barocas, Hardt, and Narayanan, \textit{Fairness and Machine Learning} (online book), selected sections.;%
Knaflic, \textit{Storytelling with Data}, selected chapters on communicating results.%
}
%
}{%
  % Faculty copies live under `private/` so their compile CWD is different.
  \IfFileExists{../../../../public/_shared/aiml-sources.tex}{%
    % Source strings for Elements of AIML (used by slides + notes).

\newcommand{\AIMLRepoURL}{https://github.com/tali7c/Elements-of-AIML}

\newcommand{\AIMLLabSources}{%
Provost and Fawcett, \textit{Data Science for Business}.;%
McKinney, \textit{Python for Data Analysis}.;%
WEKA Documentation (Waikato Environment for Knowledge Analysis), accessed 2026-02-28.;%
Matplotlib and Seaborn documentation, accessed 2026-02-28.%
}

\newcommand{\AIMLUnitISources}{%
Rich and Knight, \textit{Artificial Intelligence}, McGraw Hill, 2017.;%
Russell and Norvig, \textit{Artificial Intelligence: A Modern Approach} (4e), Ch. 1--2 (Introduction, Intelligent Agents).;%
Poole and Mackworth, \textit{Artificial Intelligence: Foundations of Computational Agents} (2e), selected sections.%
}

\newcommand{\AIMLUnitIISources}{%
Russell and Norvig, \textit{Artificial Intelligence: A Modern Approach} (4e), Ch. 7--9 (Logical Agents, First-Order Logic, Inference).;%
Huth and Ryan, \textit{Logic in Computer Science} (2e), selected sections on propositional and first-order logic.;%
Genesereth and Nilsson, \textit{Logical Foundations of Artificial Intelligence}, selected chapters.%
}

\newcommand{\AIMLUnitIIISources}{%
Mueller and Massaron, \textit{Machine Learning for Dummies}, For Dummies, 2016.;%
Mitchell, \textit{Machine Learning}, selected chapters on learning basics and evaluation.;%
Scikit-learn documentation, accessed 2026-02-28.%
}

\newcommand{\AIMLUnitIVSources}{%
Mueller and Massaron, \textit{Machine Learning for Dummies}, For Dummies, 2016.;%
James, Witten, Hastie, and Tibshirani, \textit{An Introduction to Statistical Learning} (2e), selected chapters.;%
Scikit-learn documentation, accessed 2026-02-28.%
}

\newcommand{\AIMLUnitVSources}{%
NIST, \textit{AI Risk Management Framework (AI RMF 1.0)}, 2023.;%
Barocas, Hardt, and Narayanan, \textit{Fairness and Machine Learning} (online book), selected sections.;%
Knaflic, \textit{Storytelling with Data}, selected chapters on communicating results.%
}
%
  }{%
    % Source strings for Elements of AIML (used by slides + notes).

\newcommand{\AIMLRepoURL}{https://github.com/tali7c/Elements-of-AIML}

\newcommand{\AIMLLabSources}{%
Provost and Fawcett, \textit{Data Science for Business}.;%
McKinney, \textit{Python for Data Analysis}.;%
WEKA Documentation (Waikato Environment for Knowledge Analysis), accessed 2026-02-28.;%
Matplotlib and Seaborn documentation, accessed 2026-02-28.%
}

\newcommand{\AIMLUnitISources}{%
Rich and Knight, \textit{Artificial Intelligence}, McGraw Hill, 2017.;%
Russell and Norvig, \textit{Artificial Intelligence: A Modern Approach} (4e), Ch. 1--2 (Introduction, Intelligent Agents).;%
Poole and Mackworth, \textit{Artificial Intelligence: Foundations of Computational Agents} (2e), selected sections.%
}

\newcommand{\AIMLUnitIISources}{%
Russell and Norvig, \textit{Artificial Intelligence: A Modern Approach} (4e), Ch. 7--9 (Logical Agents, First-Order Logic, Inference).;%
Huth and Ryan, \textit{Logic in Computer Science} (2e), selected sections on propositional and first-order logic.;%
Genesereth and Nilsson, \textit{Logical Foundations of Artificial Intelligence}, selected chapters.%
}

\newcommand{\AIMLUnitIIISources}{%
Mueller and Massaron, \textit{Machine Learning for Dummies}, For Dummies, 2016.;%
Mitchell, \textit{Machine Learning}, selected chapters on learning basics and evaluation.;%
Scikit-learn documentation, accessed 2026-02-28.%
}

\newcommand{\AIMLUnitIVSources}{%
Mueller and Massaron, \textit{Machine Learning for Dummies}, For Dummies, 2016.;%
James, Witten, Hastie, and Tibshirani, \textit{An Introduction to Statistical Learning} (2e), selected chapters.;%
Scikit-learn documentation, accessed 2026-02-28.%
}

\newcommand{\AIMLUnitVSources}{%
NIST, \textit{AI Risk Management Framework (AI RMF 1.0)}, 2023.;%
Barocas, Hardt, and Narayanan, \textit{Fairness and Machine Learning} (online book), selected sections.;%
Knaflic, \textit{Storytelling with Data}, selected chapters on communicating results.%
}
%
  }%
}


\title{Elements of AIML\\Unit 03 -- Lecture 01 Notes\\Introduction to Machine Learning \& Datasets}
\author{Tofik Ali}
\date{\today}

\begin{document}
\maketitle
\tableofcontents

\section{Lecture Overview}
This lecture introduces \textbf{machine learning} and the basic dataset concepts needed for the rest of the course and the labs.
We focus on: tasks, datasets, features/labels, and common data quality issues.
\notesources{\AIMLUnitIIISources}

\section{How to use these notes (self-study)}
When you study ML, always keep the full pipeline in mind:
problem definition $\rightarrow$ dataset $\rightarrow$ features/labels $\rightarrow$ training $\rightarrow$ evaluation.
Most beginner mistakes happen because one of these steps is skipped or assumed.
Try the practice questions after each major section.

\section{Learning Outcomes}
\begin{itemize}[leftmargin=*]
  \item Define ML and list key task categories.
  \item Explain datasets in terms of rows/examples and columns/variables.
  \item Distinguish features from labels and identify common feature data types.
  \item Recognize common dataset issues and why they matter.
\end{itemize}

\section{Keywords}
\begin{itemize}[leftmargin=*]
  \item \textbf{Example/instance:} one row (one data point).
  \item \textbf{Feature:} input variable used for prediction (column in $X$).
  \item \textbf{Label/target:} output variable to predict (vector $y$).
  \item \textbf{Supervised learning:} learn from labeled examples.
  \item \textbf{Unsupervised learning:} find structure without labels.
  \item \textbf{Generalization:} performance on unseen but similar data.
\end{itemize}

\section{Abbreviations and symbols}
\begin{itemize}[leftmargin=*]
  \item \textbf{ML}: machine learning.
  \item \textbf{EDA}: exploratory data analysis.
  \item $X$: feature matrix (inputs), usually written as $X \in \mathbb{R}^{n \times d}$ after encoding.
  \item $y$: target/label vector (outputs), e.g., $y \in \mathbb{R}^n$ (regression) or $y \in \{1,\dots,K\}^n$ (classification).
  \item $n$: number of examples (rows), $d$: number of features (columns), $K$: number of classes.
  \item $\hat{y}$: prediction produced by a model, $f(\cdot)$.
\end{itemize}

\section{What is machine learning?}
ML builds models from data. A common view is:
\begin{itemize}[leftmargin=*]
  \item you provide examples,
  \item you choose a learning method (algorithm),
  \item you train a model to minimize some error,
  \item you evaluate on unseen data.
\end{itemize}

\subsection{A standard definition (Mitchell)}
A classic textbook definition:
\begin{quote}
  A computer program is said to learn from \textbf{experience} $E$ with respect to some \textbf{task} $T$ and \textbf{performance measure} $P$ if its performance at task $T$, as measured by $P$, improves with experience $E$.
\end{quote}

Example:
\begin{itemize}[leftmargin=*]
  \item $T$: classify emails as spam/ham
  \item $E$: labeled emails from the past
  \item $P$: accuracy/F1 on new (unseen) emails
\end{itemize}

\subsection{Algorithm vs model (important distinction)}
\begin{itemize}[leftmargin=*]
  \item A \textbf{learning algorithm} is the procedure we run (e.g., ``fit linear regression'').
  \item A \textbf{model} is what we get after training (e.g., the learned weights $w$).
\end{itemize}
During \textbf{training}, the model parameters are estimated from data.
During \textbf{inference/prediction}, the model is used to produce $\hat{y}$ for a new input $x$.

\subsection{Why not just write rules?}
Rules are often:
\begin{itemize}[leftmargin=*]
  \item too complex to specify (vision, language),
  \item brittle when the world changes,
  \item incomplete in edge cases.
\end{itemize}
ML can learn patterns directly from data, but it introduces new responsibilities:
\begin{itemize}[leftmargin=*]
  \item collecting and cleaning datasets,
  \item choosing appropriate evaluation (so we do not fool ourselves),
  \item monitoring models after deployment (data can change over time).
\end{itemize}

\section{ML task types}
\subsection{Supervised learning}
We learn from labeled examples $(x_i, y_i)$.
\begin{itemize}[leftmargin=*]
  \item \textbf{Regression:} $y$ is numeric (price, temperature).
  \item \textbf{Classification:} $y$ is a class label (spam/ham, disease/no disease).
\end{itemize}
\paragraph{Typical outputs.}
In regression, $y \in \mathbb{R}$.
In classification, $y \in \{1,\dots,K\}$ for $K$ classes (binary classification is the case $K=2$).

\subsection{Unsupervised learning}
We use only $x_i$ (no labels).
\begin{itemize}[leftmargin=*]
  \item \textbf{Clustering:} group similar examples.
  \item \textbf{Dimensionality reduction:} compress/represent data in fewer variables.
\end{itemize}

\subsection{Other task families (brief preview)}
\begin{itemize}[leftmargin=*]
  \item \textbf{Semi-supervised learning:} a small labeled set + a larger unlabeled set.
  \item \textbf{Anomaly detection:} find rare/abnormal points (fraud, faults).
  \item \textbf{Ranking/recommendation:} order items for a user (feeds, products).
\end{itemize}

\subsection{Reinforcement learning (preview)}
An agent interacts with an environment and learns from rewards. We cover an overview later in Unit IV.

\section{Dataset basics}
\subsection{Tabular representation}
Many problems can be represented as a table:
\begin{itemize}[leftmargin=*]
  \item rows = examples/instances,
  \item columns = features + (optional) label column.
\end{itemize}
We typically denote:
\[
  X \in \mathbb{R}^{n \times d},\quad y \in \mathbb{R}^{n} \text{ or } y \in \{1,\dots,K\}^{n}
\]
where $n$ is the number of examples and $d$ is the number of features.

\subsection{What does ``$X \in \mathbb{R}^{n \times d}$'' really mean?}
Even if your raw dataset has text/categorical columns, most ML models require numeric inputs.
So we perform \textbf{encoding} (e.g., one-hot for categories, TF-IDF for text) to convert raw values into numerical feature vectors.
After encoding, we can treat $X$ as a real-valued matrix.

\subsection{Data types}
Features can be:
\begin{itemize}[leftmargin=*]
  \item \textbf{numerical} (int/float),
  \item \textbf{categorical} (city, category),
  \item \textbf{ordinal} (ordered categories),
  \item \textbf{text/images/time series} (need a feature representation).
\end{itemize}

\subsection{Dataset quality checklist (what can go wrong?)}
Common issues that strongly affect ML results:
\begin{itemize}[leftmargin=*]
  \item \textbf{Missing values:} empty/NaN cells; need imputation or removal.
  \item \textbf{Outliers:} extreme values; may be errors or rare but real cases.
  \item \textbf{Duplicates:} repeated rows can bias training and evaluation.
  \item \textbf{Label noise:} incorrect target values (wrong ground truth).
  \item \textbf{Class imbalance:} one class dominates, making accuracy misleading.
  \item \textbf{Data leakage:} information that would not exist at prediction time gets included.
\end{itemize}
Most lab experiments (EDA, missing values, outliers, scaling, imbalance) are designed to give hands-on practice with these.

\section{Feature sets and representation}
\subsection{Why representation matters}
Two models can differ more due to \textbf{features} than due to the choice of algorithm.
For example, a ``location'' feature may need \textbf{one-hot encoding} before it can be used by many models.

\subsection{Common feature transformations}
\begin{itemize}[leftmargin=*]
  \item \textbf{One-hot encoding}: convert a categorical value into a 0/1 vector.
  \item \textbf{Scaling/standardization}: make numerical features comparable (mean 0, std 1).
  \item \textbf{Log transforms}: handle skewed distributions (income, counts).
  \item \textbf{Binning}: convert continuous features into buckets (sometimes useful for interpretation).
\end{itemize}

\subsection{Preprocessing mindset (ties to labs)}
In the labs you will practice:
\begin{itemize}[leftmargin=*]
  \item importing/exporting data (Pandas),
  \item basic summaries and EDA,
  \item missing values and outlier handling,
  \item scaling and encoding.
\end{itemize}

\section{Worked example: house price prediction}
Suppose we want to predict house price:
\begin{itemize}[leftmargin=*]
  \item target $y$ = price,
  \item features $X$ = area, rooms, location, age.
\end{itemize}
If ``location'' is categorical, we cannot directly treat it as a number (that would add fake ordering). We encode it using one-hot encoding (lab experiment on encoding).

\subsection{Worked example: spam classification (conceptual)}
\begin{itemize}[leftmargin=*]
  \item Task: classify email as \{spam, ham\}.
  \item Raw input: text; label: spam/ham.
  \item Features: bag-of-words counts, TF-IDF vectors, or embeddings.
  \item Key challenge: new spam patterns appear over time (data drift), so evaluation and monitoring matter.
\end{itemize}

\section{A minimal Pandas inspection snippet (illustrative)}
\begin{lstlisting}[language=Python]
import pandas as pd

df = pd.read_csv("data.csv")
print(df.shape)        # (rows, columns)
print(df.head())       # first 5 rows
print(df.isna().sum()) # missing values per column
\end{lstlisting}
This basic habit catches many problems early.

\section{Common pitfalls (and how to avoid them)}
\begin{itemize}[leftmargin=*]
  \item \textbf{Treating categories as numbers:} encoding city as 1/2/3 creates a fake ordering; use one-hot (or other encodings).
  \item \textbf{Ignoring leakage:} features that contain future information make evaluation meaningless.
  \item \textbf{Confusing model and algorithm:} training procedure is the algorithm; the trained parameters are the model.
  \item \textbf{Assuming more features always helps:} irrelevant/noisy features can hurt generalization (dimensionality issues).
\end{itemize}

\section{Practice (with answers)}
\subsection*{P1. A dataset has 1,000 rows and 25 columns. If 1 column is the label, what is $n$ and $d$?}
\textbf{Answer:} $n=1000$. If there are 25 total columns and 1 is the label, then $d=24$ features.

\subsection*{P2. Give one example of a categorical feature and how you would encode it.}
\textbf{Answer:} Feature: city=\{Delhi, Dehradun, Mumbai\}. Encode using one-hot: three binary columns, one per city.

\subsection*{P3. Why can accuracy be misleading on imbalanced data?}
\textbf{Answer:} If 95\% examples are class 0, a model that always predicts 0 gets 95\% accuracy but is useless for detecting class 1.

\section{Exit questions (with answers)}

\subsection*{Q1. Define ML in your own words.}
\textbf{Answer:} ML is a set of methods that learn patterns from data to make predictions or decisions, rather than being explicitly programmed with rules.

\subsection*{Q2. Features vs labels?}
\textbf{Answer:} Features are input variables ($X$). Labels/targets are outputs ($y$) that the model predicts.

\subsection*{Q3. Give two supervised and two unsupervised tasks.}
\textbf{Answer:} Supervised: spam classification, house price regression. Unsupervised: customer clustering, PCA for dimensionality reduction.

\subsection*{Q4. Name three common dataset issues.}
\textbf{Answer:} missing values, outliers, noisy/incorrect labels, duplicates, inconsistent formatting, class imbalance.

\section{Summary}
ML depends on datasets and good feature representations. The next lecture focuses on splitting and evaluation to avoid misleading results.
\notesources{\AIMLUnitIIISources}

\end{document}
